\Article{22}
Die Variation der ganzen Oberfläche $U$ erhalten wir, indem wir diesen Ausdruck über alle Elemente $\mathrm{d}U$ integrieren. Zu diesem Zweck behandeln wir die beiden Teile des Integrals getrennt, nämlich
\[
\int\mathrm{d}U.((\eta\eta+\zeta\zeta)\frac{\mathrm{d}\delta x}{\mathrm{d}x}-\xi\eta\frac{\mathrm{d}\delta y}{\mathrm{d}x}-\xi\zeta\frac{\mathrm{d}\delta z}{\mathrm{d}x})=A
\]
und
\[
\int\mathrm{d}U(-\xi\eta\frac{\mathrm{d}\delta x}{\mathrm{d}y}+(\xi\xi+\zeta\zeta)\frac{\mathrm{d}\delta y}{\mathrm{d}y}-\eta\zeta\frac{\mathrm{d}\delta z}{\mathrm{d}y})=B
\]

Man denke sich eine zur Koordinatenachse $y$ senkrechte Ebene, deren fester Wert von $y$ im Bereich zwischen den Extremwerten von $y$ auf der Oberfläche $U$ liegt. Diese Ebene schneidet die Begrenzungslinie $P$ in zwei, vier, sechs usw. Punkten, deren erste Koordinaten der Reihe nach $x^0,x',x''$ usw. seien. Die übrigen zu diesen Punkten gehörigen Größen unterscheide man durch entsprechende Indizes. Ebenso werde die Oberfläche von einer anderen Ebene geschnitten, die der ersten unendlich nahe und zu ihr parallel ist und deren zweite Koordinate $y+\mathrm{d}y$ lautet. Zwischen diesen Ebenen liegen die Elemente $\mathrm{d}P^0$, $\mathrm{d}P'$, $\mathrm{d}P''$ usw. der Begrenzungslinie. Wie man leicht sieht, gilt
\[
\mathrm{d}y=-Y^0\mathrm{d}P^0=+Y'\mathrm{d}P'=-Y''\mathrm{d}P''=+Y'''\mathrm{d}P'''\quad\text{etc.}
\]
\SourceLeaf{65}Denken wir uns ferner unendlich viele zur Koordinatenachse $x$ senkrechte Ebenen, so entspricht jedem Element $\mathrm{d}x$ zwischen $x^0$ und $x'$ oder zwischen $x''$ und $x'''$ usw. ein Element $\mathrm{d}U=\frac{\mathrm{d}x.\mathrm{d}y}{\zeta}$. Daraus geht hervor, dass sich der Teil des Integrals $A$, der dem zwischen den Ebenen $y,y+\mathrm{d}y$ liegenden Teil der Oberfläche entspricht, durch die Integration
\[
\mathrm{d}y\int\mathrm{d}x.(\frac{\eta\eta+\zeta\zeta}{\zeta}.\frac{\mathrm{d}\delta x}{\mathrm{d}x}-\frac{\xi\eta}{\zeta}.\frac{\mathrm{d}\delta y}{\mathrm{d}x}-\xi\frac{\mathrm{d}\delta z}{\mathrm{d}x})
\]
ergibt, die von $x=x^0$ bis $x=x'$ und dann von $x=x''$ bis $x=x'''$ usw. erstreckt wird. Das unbestimmte Integral lautet
\[
\begin{aligned}
&(\frac{\eta\eta+\zeta\zeta}{\zeta}.\delta x-\frac{\xi\eta}{\zeta}.\delta y-\xi\delta z)\mathrm{d}y\\
&-\mathrm{d}y\int(\delta x.\frac{\mathrm{d}\frac{\eta\eta+\zeta\zeta}{\zeta}}{\mathrm{d}x}-\delta y.\frac{\mathrm{d}\frac{\xi\eta}{\zeta}}{\mathrm{d}x}-\delta z.\frac{\mathrm{d}\xi}{\mathrm{d}x})\mathrm{d}x
\end{aligned}
\]
woraus sich für unseren Fall
\[
\begin{aligned}
&(\frac{\eta^0\eta^0+\zeta^0\zeta^0}{\zeta^0}.\delta x^0-\frac{\xi^0\eta^0}{\zeta^0}.\delta y^0-\xi^0\delta z^0)Y^0\mathrm{d}P^0\\
+&(\frac{\eta'\eta'+\zeta'\zeta'}{\zeta'}.\delta x'-\frac{\xi'\eta'}{\zeta'}.\delta y'-\xi'\delta z')Y'\mathrm{d}P'\\
+&(\frac{\eta''\eta''+\zeta''\zeta''}{\zeta''}.\delta x''-\frac{\xi''\eta''}{\zeta''}.\delta y''-\xi''\delta z'')Y''\mathrm{d}P''\\
+&\text{etc.}\\
-&\int\zeta\mathrm{d}U.(\delta x.\frac{\mathrm{d}\frac{\eta\eta+\zeta\zeta}{\zeta}}{\mathrm{d}x}-\delta y.\frac{\mathrm{d}\frac{\xi\eta}{\zeta}}{\mathrm{d}x}-\delta z.\frac{\mathrm{d}\xi}{\mathrm{d}x})
\end{aligned}
\]
oder, was dasselbe ist,
\[
\begin{aligned}
&\Sigma(\frac{\eta\eta+\zeta\zeta}{\zeta}.\delta x-\frac{\xi\eta}{\zeta}.\delta y-\xi\delta z)Y\mathrm{d}P\\
&-\int\zeta\mathrm{d}U.(\delta x.\frac{\mathrm{d}\frac{\eta\eta+\zeta\zeta}{\zeta}}{\mathrm{d}x}-\delta y.\frac{\mathrm{d}\frac{\xi\eta}{\zeta}}{\mathrm{d}x}-\delta z.\frac{\mathrm{d}\xi}{\mathrm{d}x})
\end{aligned}
\]
ergibt. Dabei ist die Summe über alle Elemente $\mathrm{d}P$ und das Integral über alle Elemente $\mathrm{d}U$ zwischen den Ebenen $y$ und $y+\mathrm{d}y$ zu erstrecken.

Die gesamte Größe $A$ wird also durch
\[
\begin{aligned}
&\int(\frac{\eta\eta+\zeta\zeta}{\zeta}.\delta x-\frac{\xi\eta}{\zeta}.\delta y-\xi\delta z)Y\mathrm{d}P\\
&-\int\zeta\mathrm{d}U(\delta x.\frac{\mathrm{d}\frac{\eta\eta+\zeta\zeta}{\zeta}}{\mathrm{d}x}-\delta y.\frac{\mathrm{d}\frac{\xi\eta}{\zeta}}{\mathrm{d}x}-\delta z.\frac{\mathrm{d}\xi}{\mathrm{d}x})
\end{aligned}
\]
ausgedrückt, wobei die erste Integration über die ganze Begrenzungslinie $P$, die zweite über die ganze Oberfläche $U$ zu erstrecken ist.

\SourceLeaf{66}\Article{23}
Durch ganz entsprechende Überlegungen finden wir
\[
\begin{aligned}
B={}&\int(\frac{\xi\eta}{\zeta}.\delta x-\frac{\xi\xi+\zeta\zeta}{\zeta}.\delta y+\eta\delta z)X\mathrm{d}P\\
&+\int\zeta\mathrm{d}U(\delta x.\frac{\mathrm{d}\frac{\xi\eta}{\zeta}}{\mathrm{d}y}-\delta y.\frac{\mathrm{d}\frac{\xi\xi+\zeta\zeta}{\zeta}}{\mathrm{d}y}+\delta z.\frac{\mathrm{d}\eta}{\mathrm{d}y})
\end{aligned}
\]
Setzen wir also für einen beliebigen Punkt der Begrenzungslinie $P$
\[
\begin{aligned}
&[X\xi\eta+Y(\eta\eta+\zeta\zeta)]\delta x-[X(\xi\xi+\zeta\zeta)+Y\xi\eta]\delta y\\
&+(X\eta\zeta-Y\xi\zeta)\delta z=\zeta Q
\end{aligned}
\]
und für einen beliebigen Punkt der Oberfläche $U$
\[
\begin{aligned}
&(\frac{\mathrm{d}\frac{\xi\eta}{\zeta}}{\mathrm{d}y}-\frac{\mathrm{d}\frac{\eta\eta+\zeta\zeta}{\zeta}}{\mathrm{d}x})\zeta\delta x+(\frac{\mathrm{d}\frac{\xi\eta}{\zeta}}{\mathrm{d}x}-\frac{\mathrm{d}\frac{\xi\xi+\zeta\zeta}{\zeta}}{\mathrm{d}y})\zeta\delta y\\
&+(\frac{\mathrm{d}\xi}{\mathrm{d}x}+\frac{\mathrm{d}\eta}{\mathrm{d}y})\zeta\delta z=V
\end{aligned}
\]
so ergibt sich schließlich
\[
\delta U=\int Q\mathrm{d}P+\int V\mathrm{d}U
\]
wobei die erste Integration über die ganze Begrenzungslinie $P$ und die zweite über die ganze Oberfläche $U$ zu erstrecken ist.

\Article{24}
Die eben angegebenen Formeln für $Q$ und $V$ lassen sich erheblich verkürzen. Mithilfe der Gleichung $X\xi+Y\eta+Z\zeta=0$ nimmt $Q$ sofort die folgende symmetrische Form an.
\[
Q=(Y\zeta-Z\eta)\delta x+(Z\xi-X\zeta)\delta y+(X\eta-Y\xi)\delta z
\]

Um auch den für $V$ erhaltenen Ausdruck in eine übersichtlichere Form zu bringen, bemerken wir, dass aus den Formeln
\[
\frac{\mathrm{d}z}{\mathrm{d}x}=-\frac{\xi}{\zeta},\qquad \frac{\mathrm{d}z}{\mathrm{d}y}=-\frac{\eta}{\zeta}
\]
folgt
\[
\frac{\mathrm{d}\frac{\xi}{\zeta}}{\mathrm{d}y}=\frac{\mathrm{d}\frac{\eta}{\zeta}}{\mathrm{d}x}
\]
Daraus ergibt sich
\[
\frac{\mathrm{d}\frac{\xi\eta}{\zeta}}{\mathrm{d}y}=\frac{\xi}{\zeta}.\frac{\mathrm{d}\eta}{\mathrm{d}y}+\eta\frac{\mathrm{d}\frac{\xi}{\zeta}}{\mathrm{d}y}=\frac{\xi}{\zeta}.\frac{\mathrm{d}\eta}{\mathrm{d}y}+\eta\frac{\mathrm{d}\frac{\eta}{\zeta}}{\mathrm{d}x}
\]
\SourceLeaf{67}Ferner leiten wir aus $\xi\xi+\eta\eta+\zeta\zeta=1$ ab
\[
\xi\frac{\mathrm{d}\xi}{\mathrm{d}x}+\eta\frac{\mathrm{d}\eta}{\mathrm{d}x}+\zeta\frac{\mathrm{d}\zeta}{\mathrm{d}x}=0
\]
und daraus
\[
\frac{\mathrm{d}\frac{\eta\eta+\zeta\zeta}{\zeta}}{\mathrm{d}x}=\eta\frac{\mathrm{d}\frac{\eta}{\zeta}}{\mathrm{d}x}+\frac{\eta}{\zeta}.\frac{\mathrm{d}\eta}{\mathrm{d}x}+\frac{\mathrm{d}\zeta}{\mathrm{d}x}=\eta\frac{\mathrm{d}\frac{\eta}{\zeta}}{\mathrm{d}x}-\frac{\xi}{\zeta}.\frac{\mathrm{d}\xi}{\mathrm{d}x}
\]
Setzen wir diese Werte in den Koeffizienten von $\delta x$ im Ausdruck für $V$ ein, so ergibt sich für ihn
\[
=\xi(\frac{\mathrm{d}\xi}{\mathrm{d}x}+\frac{\mathrm{d}\eta}{\mathrm{d}y})
\]
Ganz entsprechend geht der Koeffizient von $\delta y$ in demselben Ausdruck über in
\[
\eta(\frac{\mathrm{d}\xi}{\mathrm{d}x}+\frac{\mathrm{d}\eta}{\mathrm{d}y})
\]
Damit erhalten wir
\[
V=(\xi\delta x+\eta\delta y+\zeta\delta z)(\frac{\mathrm{d}\xi}{\mathrm{d}x}+\frac{\mathrm{d}\eta}{\mathrm{d}y})
\]

